\clearpage
\renewcommand{\sectioncolor}{alert}
\section{Where mathematics provably cannot help}
\markboth{The repeat wall}{}

This section exists so that nobody on the team spends six months on a problem that has been proved
unsolvable. Two of our limits are not engineering limits. They are mathematical facts.

\begin{alertbox}[title={Limit 1 --- repeats longer than the read cannot be resolved. At all.}]
Bresler, Bresler \& Tse (\textit{BMC Bioinformatics}, 2013) proved the feasibility condition for
shotgun assembly: reconstruction is possible \emph{only if} the read length exceeds the longest
interleaved and triple repeats.
\begin{center}
\begin{tikzpicture}[font=\footnotesize]
 \draw[slate!50,line width=1pt,-{Stealth[length=5pt]}] (0,0)--(14.6,0);
 \node[font=\scriptsize,text=slate,anchor=north] at (7.3,-0.75) {length (bp, log scale)};
 \foreach \x/\l in {1.0/{100},4.2/{500},7.4/{2{,}000},10.6/{10{,}000},13.4/{50{,}000}}{
   \draw[slate!50] (\x,-0.12)--(\x,0.12); \node[font=\scriptsize,below=3pt] at (\x,-0.1) {\l};}
 \fill[bio] (2.4,0) circle (2.6pt);
 \node[font=\scriptsize\bfseries,text=bio,anchor=south] at (2.4,0.22) {our reads: 233 bp};
 \fill[alert] (7.7,0) circle (2.6pt);
 \node[font=\scriptsize\bfseries,text=alert,anchor=south] at (7.7,0.22) {NODE\_18: 2{,}600 bp $\times$ 8};
 \fill[alert] (11.6,0) circle (2.6pt);
 \node[font=\scriptsize\bfseries,text=alert,anchor=south] at (11.6,0.22) {NODE\_14: 19{,}885 bp $\times$ 3};
 \draw[alert,line width=1.2pt,decorate,decoration={brace,amplitude=5pt,mirror}] (2.55,-0.35)--(11.45,-0.35);
 \node[font=\scriptsize\bfseries,text=alert,anchor=north] at (7.0,-0.55) {};
\end{tikzpicture}
\end{center}
\vspace{-6pt}
\textbf{No algorithm, no AI model, and no amount of coverage closes that gap.} If we need a closed
assembled chromosome, we need \emph{longer reads} --- roughly 20$\times$ of nanopore or PacBio, hybrid
assembled with what we already have. That is a procurement decision, not a research problem.
\end{alertbox}

\begin{alertbox}[title={Limit 2 --- photon shot noise}]
From our own handbook, Chapter T6.2: for a cluster yielding $N$ photoelectrons,
$\mathrm{SNR}\approx\sqrt{N}$. \textbf{To double the signal-to-noise ratio you need four times the
photons.} This is a property of light, not of Illumina. Clusters (about 1{,}000 copies) already buy
a factor of $\sqrt{1000}\approx32$. There is no further mathematics to extract here.
\end{alertbox}

\begin{plain}
Two of our problems are worth working on and two are not. Spending effort on cleverer algorithms to
resolve the rRNA operons from 233 bp reads is like trying to read the small print in a photograph
that was taken out of focus --- the information is not in the file. Buy a long-read run instead.
\end{plain}
